\scanpage{30}
\mspage{11}

To construct any Quillen triple, we may proceed as follows: Start with \inserted{a localizer} $C\subset\mathrm{Fl}(\mathscr A)$, \textcircled{a} \deleted{\unsure{being}}\inserted{$C$ stable} \deleted{\unsure{$C$}} by direct factors, hence a pair $(C,C_*=TF)$ --- we \uline{must check \textcircled{b} factorization} \uline{condition for the pair} $(C,C_*)$. \deleted{\unsure{Now|Hence}} Thus $(C,TF)$ is a Quillen pair.

Let next \textcircled{c} $\boxed{W\supset TF}$ any localizer which is \inserted{\textcircled{d} \uline{$W$ mildly}} \uline{saturated and stable under direct factor}. We now get $TC=C\cap W$ \inserted{$\subset C$}, which is \deleted{\unsure{again}} stable under direct factor. Take $F=(TC)_*$ \inserted{$(\supset TF=C_*)$}, we must check \textcircled{\unsure{e|d}} \uline{factorization property for the pair} $(TC,F)$ --- thus we get that $(TC,F)$ is a Quillen pair. Lastly, we must still check that $TF=F\cap W$, and by c) we know already $TF\subset F\cap W$, we have to check still that $F\cap W\subset TF$. But using factorization $f=pi,\ p\in TF,\ i\in C$, as $f,p\in W$ we get $i\in W$ hence $i\in W\cap C=TC$, \deleted{hence} there exists $s:X'-X$ s that $si=\mathrm{id}_X$ (retraction of $X'$ upon $X$) and $fs=p$, hence $f$ a retraction of $p$, hence is in $TF$, OK. Thus

\diagram{\begin{tikzcd}[ampersand replacement=\&,row sep=small,column sep=large]
X\arrow[r,dash,"="]\arrow[d,"{i\in TC}"'] \& X\arrow[d,"{f\in F}"] \& \& X'\arrow[r,shift left,"s"]\arrow[d,dash,"p"'] \& X\arrow[l,bend left=20,"i"]\arrow[d,dash,"f"]\\
X'\arrow[r,"p"']\arrow[ur,dashed,"s"]\&Y \&\&Y\arrow[r,shift left,"="]\&Y\arrow[l,shift left]
\end{tikzcd}}

\uline{Proposition} In order for a pair \inserted{of localizers} $(C,W)$ to correspond to a Quillen triple $(C,F,W)$, it is ns that the following conditions be satisfied
\[
\left\{\begin{aligned}
1)&\ C,W\ \text{\uline{stable by direct factors}},\ \deleted{\unsure{+}}\ W\text{ mildly saturated}\\
2)&\ \text{Factorization condition for }(C,C_*)\text{ and for }(C\cap W,(C\cap W)_*).\\
3)&\ W\supset TF\overset{\mathrm{def}}=C_*
\end{aligned}\right.
\]

We have to develop mainly a criterion to get factorization. [In the case of interest to us, we get 1) free, 3) almost free (there comes in the ``test''-assumption though), and 2) causes a little technical problem. I am however more than in factorization conditions, interested in the
% archive sheet number: 29
